\section{Discussion}

The change with the widest consequences here is a policy, not a number: no
thermodynamic source is promoted over another. The 2020 release published one
accepted value per reaction; this one publishes four, each with its own
uncertainty and direction call, because the sources disagree in ways a single
value conceals --- and because the promotion rule rewarded whichever source
reported the smallest error, which is the one that understates it.

Grading against experiment makes that tractable. A consumer wanting one number
takes the highest-graded source and knows what the grade means; one propagating
uncertainty has the covariance and the caveats. It costs little and prevents the
common misuse of treating a group-contribution estimate and an anchored
eQuilibrator value as interchangeable.

Six years of use suggest the integration role is the one the community
exercises. The 2020 release has been cited 193 times at an undiminished rate,
and among citing papers with open full text 38\% use the database in their
methods, overwhelmingly for genome-scale reconstruction. GEMsembler builds
consensus models across tools, exploiting the fact that independent pipelines
share the ModelSEED namespace to discard ambiguous identifier
mappings~\cite{gemsembler2025} --- the ``Rosetta Stone'' function, exercised by
someone other than us. Yeast9 assigned Gibbs energies to 98\% of its
metabolites and 97\% of its reactions, taking from this database the metabolite
values its other sources lacked~\cite{yeast9}. That second case is the rarer
one, and the thermodynamic layer being used by a minority of citing work is
part of why this release concentrates on making it auditable rather than
larger.

Two limitations should be read alongside the results. The protonation layer is
no longer open-source at all. Every ladder behind the shipped energies comes
from one licensed tool, and that is a deliberate exchange: ladders drawn from
different models do not cancel across a reaction, so mixing them to raise
coverage introduces an error that no single compound's record reveals. We chose
internal consistency and disclosed the cost rather than the reverse. The choice
is auditable and reversible --- the open predictions ship for every compound
with a structure --- but closing the gap properly needs measured macroscopic
ladders for compounds no published compilation covers, the nicotinamide
cofactors chief among them. And 25,855 reactions still receive no direction from
any source. Neither gap is closed by better prediction: both need measurements
or curation.

The database is a substrate rather than a product, and this version is designed
to make the uncertainty in it visible rather than tidy.
